
\begin{longtable}[]{@{}
  >{\raggedright\arraybackslash}p{(\linewidth - 2\tabcolsep) * \real{0.5000}}
  >{\raggedright\arraybackslash}p{(\linewidth - 2\tabcolsep) * \real{0.5000}}@{}}
\toprule\noalign{}
\endhead
\bottomrule\noalign{}
\endlastfoot
\multicolumn{2}{@{}>{\raggedright\arraybackslash}p{(\linewidth - 2\tabcolsep) * \real{1.0000} + 2\tabcolsep}@{}}{%
\begin{minipage}[t]{\linewidth}\raggedright
波恩哈德·黎曼（Bernhard Riemann）
\end{minipage}} \\
\multicolumn{2}{@{}>{\raggedright\arraybackslash}p{(\linewidth - 2\tabcolsep) * \real{1.0000} + 2\tabcolsep}@{}}{%
\begin{minipage}[t]{\linewidth}\raggedright
{\href{https://en.wikipedia.org/wiki/File:Georg_Friedrich_Bernhard_Riemann.jpeg}{}}

黎曼，约 1863 年
\end{minipage}} \\
出生 & \begin{minipage}[t]{\linewidth}\raggedright
格奥尔格·弗里德里希·波恩哈德·黎曼

\hfill\break
1826 年 9 月 17 日\\

\href{https://en.wikipedia.org/wiki/Breselenz}{布雷瑟伦茨（Breselenz）}，汉诺威王国（今德国）

{}\strut
\end{minipage} \\
逝世 & \begin{minipage}[t]{\linewidth}\raggedright
1866 年 7 月 20 日（享年 39 岁）\\

\href{https://en.wikipedia.org/wiki/Verbania}{塞拉斯卡（Selasca）}，意大利王国

{}\strut
\end{minipage} \\
就读学校 & \begin{minipage}[t]{\linewidth}\raggedright
\begin{itemize}
\tightlist
\item
  \href{https://en.wikipedia.org/wiki/University_of_Göttingen}{哥廷根大学}
\item
  \href{https://en.wikipedia.org/wiki/University_of_Berlin}{柏林大学}
\end{itemize}
\end{minipage} \\
著名成就 & \emph{\href{https://en.wikipedia.org/wiki/List_of_topics_named_after_Bernhard_Riemann}{见列表}} \\
\multicolumn{2}{@{}>{\raggedright\arraybackslash}p{(\linewidth - 2\tabcolsep) * \real{1.0000} + 2\tabcolsep}@{}}{%
\textbf{科学生涯}} \\
研究领域 & \begin{minipage}[t]{\linewidth}\raggedright
\begin{itemize}
\tightlist
\item
  \href{https://en.wikipedia.org/wiki/Mathematics}{数学}
\item
  \href{https://en.wikipedia.org/wiki/Physics}{物理学}
\end{itemize}
\end{minipage} \\
任职机构 & \href{https://en.wikipedia.org/wiki/University_of_Göttingen}{哥廷根大学} \\
\href{https://en.wikipedia.org/wiki/Thesis}{博士论文} & \emph{\href{http://www.maths.tcd.ie/pub/HistMath/People/Riemann/Grund/Grund.pdf}{Grundlagen für eine allgemeine Theorie der Funktionen einer veränderlichen complexen Größe}（一个复变量函数一般理论的基础）}{~}{(1851)} \\
\href{https://en.wikipedia.org/wiki/Doctoral_advisor}{博士导师} & \href{https://en.wikipedia.org/wiki/Carl_Friedrich_Gauss}{卡尔·弗里德里希·高斯} \\
其他学术导师 & \begin{minipage}[t]{\linewidth}\raggedright
\begin{itemize}
\tightlist
\item
  \href{https://en.wikipedia.org/wiki/Gotthold_Eisenstein}{戈特霍尔德·艾森斯坦（Gotthold Eisenstein）}
\item
  \href{https://en.wikipedia.org/wiki/Moritz_Abraham_Stern}{莫里茨·A·斯特恩（Moritz A. Stern）}
\item
  \href{https://en.wikipedia.org/wiki/Carl_Wolfgang_Benjamin_Goldschmidt}{卡尔·W·B·戈尔德施密特（Carl W. B. Goldschmidt）}
\end{itemize}
\end{minipage} \\
著名学生 & \begin{minipage}[t]{\linewidth}\raggedright
\href{https://en.wikipedia.org/wiki/Gustav_Roch}{古斯塔夫·罗赫（Gustav Roch）}\\
\href{https://en.wikipedia.org/wiki/Eduard_Selling}{爱德华·泽林（Eduard Selling）}\strut
\end{minipage} \\
\multicolumn{2}{@{}>{\raggedright\arraybackslash}p{(\linewidth - 2\tabcolsep) * \real{1.0000} + 2\tabcolsep}@{}}{%
} \\
\multicolumn{2}{@{}>{\raggedright\arraybackslash}p{(\linewidth - 2\tabcolsep) * \real{1.0000} + 2\tabcolsep}@{}}{%
签名} \\
\multicolumn{2}{@{}>{\raggedright\arraybackslash}p{(\linewidth - 2\tabcolsep) * \real{1.0000} + 2\tabcolsep}@{}}{%
{\href{https://en.wikipedia.org/wiki/File:Bernhard_Riemann_signature.png}{\includegraphics[width=1.5625in,height=0.35417in]{250px-Bernhard_Riemann_signature.png}}}} \\
\end{longtable}

\textbf{格奥尔格·弗里德里希·波恩哈德·黎曼}（Georg Friedrich Bernhard Riemann，1826 年 9 月 17 日 -- 1866 年 7 月 20 日）是一位德国\href{https://en.wikipedia.org/wiki/Mathematician}{数学家}，在\href{https://en.wikipedia.org/wiki/Mathematical_analysis}{分析学}、\href{https://en.wikipedia.org/wiki/Number_theory}{数论}与\href{https://en.wikipedia.org/wiki/Differential_geometry}{微分几何}诸领域均做出了深远的贡献。在\href{https://en.wikipedia.org/wiki/Real_analysis}{实分析}领域，他最为人熟知的工作是首次给出了积分的严格定义——\href{https://en.wikipedia.org/wiki/Riemann_integral}{黎曼积分}，以及他对\href{https://en.wikipedia.org/wiki/Fourier_series}{傅里叶级数}的研究。在\href{https://en.wikipedia.org/wiki/Complex_analysis}{复分析}方面，他最为重要的贡献是引入了\href{https://en.wikipedia.org/wiki/Riemann_surface}{黎曼曲面}，开创性地以一种自然的、几何的方式来处理复分析。他于 \href{https://en.wikipedia.org/wiki/On_the_Number_of_Primes_Less_Than_a_Given_Magnitude}{1859 年关于\href{https://en.wikipedia.org/wiki/Prime-counting_function}{素数计数函数}的论文}最早陈述了\href{https://en.wikipedia.org/wiki/Riemann_hypothesis}{黎曼猜想}，被视为\href{https://en.wikipedia.org/wiki/Analytic_number_theory}{解析数论}的奠基之作。通过他在\href{https://en.wikipedia.org/wiki/Riemannian_Geometry}{微分几何方面的开创性贡献}，黎曼为\href{https://en.wikipedia.org/wiki/General_relativity}{广义相对论}的数学奠定了基础。许多人认为他是有史以来最伟大的数学家之一。


\subsection{早年生活}\label{Early_years}

黎曼于 1826 年 9 月 17 日出生在\href{https://en.wikipedia.org/wiki/Kingdom_of_Hanover}{汉诺威王国}\href{https://en.wikipedia.org/wiki/Dannenberg_(Elbe)}{丹嫩贝格（Dannenberg）}附近的小村庄\href{https://en.wikipedia.org/wiki/Breselenz}{布雷瑟伦茨（Breselenz）}。父亲弗里德里希·波恩哈德·黎曼是布雷瑟伦茨的一位贫穷的\href{https://en.wikipedia.org/wiki/Lutheranism}{路德宗}牧师，曾参加\href{https://en.wikipedia.org/wiki/Napoleonic_Wars}{拿破仑战争}。母亲夏洛特·埃贝尔（Charlotte Ebell）于 1846 年去世。黎曼是六个孩子中的老二。他自幼便展现出超凡的数学天赋（如惊人的计算能力），但生性害羞、害怕在公众面前讲话，且体弱多病。


\subsection{求学经历}\label{Education}

1840 年，黎曼来到\href{https://en.wikipedia.org/wiki/Hanover}{汉诺威}与祖母同住，并就读\href{https://en.wikipedia.org/wiki/Lyceum}{文理中学}（相当于初中阶段），因为这种学校在他家所在的村庄是没有的。1842 年祖母去世后，他转入\href{https://en.wikipedia.org/wiki/Lüneburg}{吕讷堡（Lüneburg）}的\href{https://de.wikipedia.org/wiki/Johanneum\%20Lüneburg}{约翰内乌姆·吕讷堡（Johanneum Lüneburg）}继续高中学业。在那里，黎曼深入研读\href{https://en.wikipedia.org/wiki/Bible}{《圣经》}，但常常被数学吸引而分心。他的老师们对他能够进行复杂数学运算感到惊叹，他在数学上的造诣经常超出老师本身的知识水平。1846 年，19 岁的他开始学习\href{https://en.wikipedia.org/wiki/Philology}{语文学}与\href{https://en.wikipedia.org/wiki/Christian_theology}{基督教神学}，准备成为牧师以补贴家用。

1846 年春，父亲攒够了钱，把黎曼送到\href{https://en.wikipedia.org/wiki/University_of_Göttingen}{哥廷根大学}，原本计划修读\href{https://en.wikipedia.org/wiki/Theology}{神学}学位。然而到了哥廷根之后，他开始在\href{https://en.wikipedia.org/wiki/Carl_Friedrich_Gauss}{卡尔·弗里德里希·高斯}门下学习\href{https://en.wikipedia.org/wiki/Mathematics}{数学}（具体地，是高斯关于\href{https://en.wikipedia.org/wiki/Method_of_least_squares}{最小二乘法}的讲座）。高斯建议黎曼放弃神学转攻数学；在征得父亲同意后，黎曼于 1847 年转学至\href{https://en.wikipedia.org/wiki/University_of_Berlin}{柏林大学}。求学期间正值\href{https://en.wikipedia.org/wiki/Carl_Gustav_Jacob_Jacobi}{卡尔·古斯塔夫·雅可比（C. G. J. Jacobi）}、\href{https://en.wikipedia.org/wiki/Peter_Gustav_Lejeune_Dirichlet}{彼得·古斯塔夫·勒热纳·狄利克雷（Dirichlet）}、\href{https://en.wikipedia.org/wiki/Jakob_Steiner}{雅各布·施泰纳（Jakob Steiner）}与\href{https://en.wikipedia.org/wiki/Gotthold_Eisenstein}{戈特霍尔德·艾森斯坦}在该校任教。他在柏林学习了两年，于 1849 年返回哥廷根。


\subsection{学术生涯}\label{Academia}

1854 年，黎曼做了他的首场学术讲座，由此开创了\href{https://en.wikipedia.org/wiki/Riemannian_geometry}{黎曼几何}这一领域，从而为\href{https://en.wikipedia.org/wiki/Albert_Einstein}{阿尔伯特·爱因斯坦}的\href{https://en.wikipedia.org/wiki/General_theory_of_relativity}{广义相对论}奠定了舞台。1857 年，校方曾试图把黎曼提升为\href{https://en.wikipedia.org/wiki/University_of_Göttingen}{哥廷根大学}的"特任教授"（extraordinary professor）。这次尝试虽未成功，但终于让黎曼获得了固定薪水。1859 年狄利克雷（其时任高斯在哥廷根的讲席）去世后，黎曼被擢升为哥廷根大学数学系系主任。他也是第一个提出使用\href{https://en.wikipedia.org/wiki/Higher_dimensions}{高于三维或四维的维数}来描述物理实在的人。

1862 年，他与伊莉泽·科赫（Elise Koch）结婚，育有一女。


\subsection{新教家庭与在意大利的辞世}\label{Protestant_family_and_death_in_Italy}

\begin{figure}
\centering
\href{https://en.wikipedia.org/wiki/File:2009_11_riemann_008.jpg}{\includegraphics[width=2.60417in,height=3.46875in]{250px-2009_11_riemann_008.jpg}}
\caption{黎曼位于意大利\href{https://en.wikipedia.org/wiki/Piedmont}{皮埃蒙特}\href{https://en.wikipedia.org/wiki/Biganzolo}{比甘佐洛（Biganzolo）}的墓碑}\label{mwdg}
\end{figure}

1866 年\href{https://en.wikipedia.org/wiki/Hanover}{汉诺威}与\href{https://en.wikipedia.org/wiki/Prussia}{普鲁士}两国军队在哥廷根交战时，黎曼避走他乡。他在第三次旅居意大利期间，因\href{https://en.wikipedia.org/wiki/Tuberculosis}{肺结核}病逝于塞拉斯卡（Selasca，今属\href{https://en.wikipedia.org/wiki/Lake_Maggiore}{马焦雷湖}畔的\href{https://en.wikipedia.org/wiki/Verbania}{韦尔巴尼亚（Verbania）}的一个小村），葬于韦尔巴尼亚的比甘佐洛公墓。\\

黎曼是一位虔诚的基督徒，是\href{https://en.wikipedia.org/wiki/Protestant}{新教}牧师之子，他将自己作为数学家的人生视作侍奉上帝的另一种方式。他终其一生秉持基督教信仰，并视之为人生最重要的部分。临终时，他正与妻子一同诵读\href{https://en.wikipedia.org/wiki/Lord's_Prayer}{《主祷文》}，未及念完便已离世。与此同时，在哥廷根，他的女管家把他办公室里的一些文稿丢弃了，其中包括许多未发表的工作。黎曼一向不愿发表不完整的工作，因此一些深刻的洞见可能就此遗失。


\subsection{黎曼几何}\label{Riemannian_geometry}

黎曼发表的著作开辟了将分析学与几何学结合的新研究领域。这些工作日后成为了\href{https://en.wikipedia.org/wiki/Riemannian_geometry}{黎曼几何}、\href{https://en.wikipedia.org/wiki/Algebraic_geometry}{代数几何}以及\href{https://en.wikipedia.org/wiki/Complex_manifold}{复流形}理论的核心组成部分。\href{https://en.wikipedia.org/wiki/Riemann_surface}{黎曼曲面}的理论后来由\href{https://en.wikipedia.org/wiki/Felix_Klein}{费利克斯·克莱因（Felix Klein）}、特别是\href{https://en.wikipedia.org/wiki/Adolf_Hurwitz}{阿道夫·胡尔维茨（Adolf Hurwitz）}进一步发展。这一数学领域是\href{https://en.wikipedia.org/wiki/Topology}{拓扑学}的基础之一，至今仍以新的方式被应用于\href{https://en.wikipedia.org/wiki/Mathematical_physics}{数学物理}。

1853 年，\href{https://en.wikipedia.org/wiki/Carl_Friedrich_Gauss}{高斯}请其学生黎曼准备一篇关于几何基础的\emph{\href{https://en.wikipedia.org/wiki/Habilitationsschrift}{特许任教资格论文（Habilitationsschrift）}}。在历经数月的工作后，黎曼发展出了关于\href{https://en.wikipedia.org/wiki/Higher_dimensions}{高维空间}的理论，并于 1854 年 6 月 10 日在哥廷根做了著名的就职讲演 \emph{Ueber die Hypothesen, welche der Geometrie zu Grunde liegen}（"论作为几何学基础的假设"）。这篇讲演直到 12 年后（1868 年）才由戴德金（Dedekind）在他逝世两年后整理出版。它最初的反响似乎不大，但今天已被公认为几何学史上最重要的著作之一。

由这篇讲演开创的学科正是\href{https://en.wikipedia.org/wiki/Riemannian_geometry}{黎曼几何}。黎曼找到了把曲面的\href{https://en.wikipedia.org/wiki/Differential_geometry}{微分几何}（高斯本人在其\emph{\href{https://en.wikipedia.org/wiki/Theorema_egregium}{绝妙定理（theorema egregium）}}中已经证明）正确推广到 \emph{n} 维的方法。其中的基本对象就是\href{https://en.wikipedia.org/wiki/Riemannian_metric}{黎曼度量（Riemannian metric）}与\href{https://en.wikipedia.org/wiki/Riemann_curvature_tensor}{黎曼曲率张量}。在曲面（二维）情形，每点的曲率可归结为一个数（标量），而具有恒正或恒负曲率的曲面正是\href{https://en.wikipedia.org/wiki/Non-Euclidean_geometry}{非欧几何}的模型。

黎曼度量是空间每一点处的一组数（即一个\href{https://en.wikipedia.org/wiki/Tensor}{张量}），它使得我们能够测量沿任意轨迹的速度，进而通过积分给出该轨迹两端点间的距离。例如，黎曼发现：在四维空间中，无论\href{https://en.wikipedia.org/wiki/Manifold}{流形}如何弯曲，都需要每点 10 个数才能描述它的距离与曲率。


\subsection{复分析}\label{Complex_analysis}

在博士论文中，他通过\href{https://en.wikipedia.org/wiki/Riemann_surface}{黎曼曲面}为\href{https://en.wikipedia.org/wiki/Complex_analysis}{复分析}建立了几何基础——借助黎曼曲面，多值函数（如\href{https://en.wikipedia.org/wiki/Logarithm}{对数函数}（无穷多叶）或\href{https://en.wikipedia.org/wiki/Square_root}{平方根}（两叶））可以化为\href{https://en.wikipedia.org/wiki/One-to-one_function}{一一映射}的函数。在这些曲面上，复函数是\href{https://en.wikipedia.org/wiki/Harmonic_functions}{调和函数}（即满足\href{https://en.wikipedia.org/wiki/Laplace's_equation}{拉普拉斯方程}，因而满足\href{https://en.wikipedia.org/wiki/Cauchy–Riemann_equations}{柯西--黎曼方程}），它们由其奇点的位置和曲面的拓扑所确定。黎曼曲面的拓扑"亏格"由 {{\(g = w/2 - n + 1\)}{}} 给出，其中曲面有 {{\(n\)}{}} 叶在 {{\(w\)}{}} 个分支点处汇合。当 {{\(g > 1\)}{}} 时，黎曼曲面具有 {{\((3g - 3)\)}{}} 个参数（即"\href{https://en.wikipedia.org/wiki/Moduli_of_algebraic_curves}{模}"）。

他在这一领域的贡献不胜枚举。著名的\href{https://en.wikipedia.org/wiki/Riemann_mapping_theorem}{黎曼映射定理}指出：复平面上任何一个单连通区域都"双全纯等价"于（即存在双向都全纯的双射至）{{\(\mathbb{C}\)}{}} 或单位圆的内部。该定理在黎曼曲面上的推广就是著名的\href{https://en.wikipedia.org/wiki/Uniformization_theorem}{单值化定理}，它在 19 世纪由\href{https://en.wikipedia.org/wiki/Henri_Poincaré}{亨利·庞加莱（Henri Poincaré）}与\href{https://en.wikipedia.org/wiki/Felix_Klein}{费利克斯·克莱因}所证明。这里同样地，严格证明只有在更丰富的数学工具（在此为拓扑学）建立之后才得以给出。在证明黎曼曲面上函数存在性时，黎曼使用了一个极小性条件，他称之为\href{https://en.wikipedia.org/wiki/Dirichlet_principle}{狄利克雷原理}。\href{https://en.wikipedia.org/wiki/Karl_Weierstrass}{卡尔·魏尔斯特拉斯（Karl Weierstrass）}发现该证明存在漏洞：黎曼没有意识到他的工作假设（极小值存在）可能不成立——所讨论的函数空间可能不完备，因此极小值的存在性并未得到保证。直到\href{https://en.wikipedia.org/wiki/David_Hilbert}{大卫·希尔伯特（David Hilbert）}在变分法上的工作出现，狄利克雷原理才最终被严格确立起来。除此之外，魏尔斯特拉斯对黎曼非常推崇，尤其欣赏他关于\href{https://en.wikipedia.org/wiki/Abelian_function}{阿贝尔函数}的理论。当黎曼的著作问世时，魏尔斯特拉斯撤回了自己投给\emph{\href{https://en.wikipedia.org/wiki/Crelle's_Journal}{《克雷尔杂志》}}的论文，未予发表。1859 年黎曼访问柏林时，两人交流甚为投契。魏尔斯特拉斯还鼓励他的学生\href{https://en.wikipedia.org/wiki/Hermann_Amandus_Schwarz}{赫尔曼·阿曼杜斯·施瓦茨（H. A. Schwarz）}寻找复分析中代替狄利克雷原理的方法，施瓦茨成功地做到了这一点。\href{https://en.wikipedia.org/wiki/Arnold_Sommerfeld}{阿诺尔德·索末菲（Arnold Sommerfeld）}讲过的一则趣事可以反映同时代数学家在理解黎曼新思想时的困难：1870 年，魏尔斯特拉斯在去里吉山度假时随身带着黎曼的博士论文，抱怨说"难以理解"。物理学家\href{https://en.wikipedia.org/wiki/Hermann_von_Helmholtz}{赫尔曼·冯·亥姆霍兹（Hermann von Helmholtz）}帮他通宵研读，第二天回来对他说：这论文"完全自然"且"非常容易理解"。

其他亮点还包括他关于阿贝尔函数与黎曼曲面上\href{https://en.wikipedia.org/wiki/Theta_functions}{theta 函数（θ 函数）}的工作。自 1857 年起，黎曼与魏尔斯特拉斯就在比拼解决阿贝尔积分（即\href{https://en.wikipedia.org/wiki/Elliptic_integrals}{椭圆积分}的推广）的雅可比反演问题。黎曼采用多变量 theta 函数，将问题化归为求这些 theta 函数的零点。他还研究了周期矩阵，并通过"黎曼周期关系"（对称且实部为负）刻画了它们。\href{https://en.wikipedia.org/wiki/Ferdinand_Georg_Frobenius}{费迪南德·格奥尔格·弗罗贝尼乌斯（Frobenius）}与\href{https://en.wikipedia.org/wiki/Solomon_Lefschetz}{所罗门·莱夫谢茨（Solomon Lefschetz）}指出：该关系成立，等价于商空间 {{\(\mathbb{C}^{n}/\Omega\)}{}}（其中 {{\(\Omega\)}{}} 是周期矩阵生成的格）能通过 theta 函数嵌入到一个射影空间中。对于某些 {{\(n\)}{}} 值，这就是黎曼曲面的\href{https://en.wikipedia.org/wiki/Jacobian_variety}{雅可比簇（Jacobian variety）}——阿贝尔流形的一个例子。

许多数学家——例如\href{https://en.wikipedia.org/wiki/Alfred_Clebsch}{阿尔弗雷德·克莱布施（Alfred Clebsch）}——延续了黎曼在代数曲线方面的工作。这些理论都依赖于在黎曼曲面上定义的函数的性质。比如\href{https://en.wikipedia.org/wiki/Riemann–Roch_theorem}{黎曼--罗赫定理}（罗赫为黎曼的学生）描述了在黎曼曲面上、对零点和极点附加已知条件的线性独立微分的个数。

按\href{https://en.wikipedia.org/wiki/Detlef_Laugwitz}{德特勒夫·劳格维茨（Detlef Laugwitz）}的说法，\href{https://en.wikipedia.org/wiki/Automorphic_function}{自守函数}首次出现于一篇关于带电圆柱面上拉普拉斯方程的论文中。然而黎曼自己曾在 1859 年关于超几何函数的讲演以及在他关于\href{https://en.wikipedia.org/wiki/Minimal_surface}{极小曲面}的论文中，把这些函数用于共形映射（例如把拓扑三角形映到圆上）。


\subsection{实分析}\label{Real_analysis}

在\href{https://en.wikipedia.org/wiki/Real_analysis}{实分析}领域，他在自己的\href{https://en.wikipedia.org/wiki/Habilitation}{特许任教资格论文}中发现了\href{https://en.wikipedia.org/wiki/Riemann_integral}{黎曼积分}。他证明了：每一个分段连续函数都是可积的。类似地，\href{https://en.wikipedia.org/wiki/Stieltjes_integral}{斯蒂尔切斯（Stieltjes）积分}也可追溯到这位哥廷根数学家，因此两者合称为\href{https://en.wikipedia.org/wiki/Riemann–Stieltjes_integral}{黎曼--斯蒂尔切斯积分}。

在他关于\href{https://en.wikipedia.org/wiki/Fourier_series}{傅里叶级数}的特许任教资格论文中——这一研究继承了他老师狄利克雷的工作——他证明了：黎曼可积的函数都能"由"傅里叶级数"表示"。狄利克雷此前已对连续且分段可微的函数（即仅有可数多个不可微点）证明了这一点；黎曼则给出了一个例子：傅里叶级数表示一个连续但几乎处处不可微的函数——这是狄利克雷未曾涵盖的情形。他还证明了\href{https://en.wikipedia.org/wiki/Riemann–Lebesgue_lemma}{黎曼--勒贝格引理}：如果一个函数能由傅里叶级数表示，那么当 \emph{n} 很大时，其傅里叶系数趋向于零。

黎曼的这篇论文也成为\href{https://en.wikipedia.org/wiki/Georg_Cantor}{格奥尔格·康托尔（Georg Cantor）}研究傅里叶级数的起点，并最终促成\href{https://en.wikipedia.org/wiki/Set_theory}{集合论}的诞生。

1857 年，他还使用复分析方法研究了\href{https://en.wikipedia.org/wiki/Hypergeometric_differential_equations}{超几何微分方程}，并通过环绕奇点的闭路径行为（用\href{https://en.wikipedia.org/wiki/Monodromy_matrix}{单值化矩阵（monodromy matrix）}描述）给出方程的解。"由先验给定的单值化矩阵反求微分方程的存在性问题"正是希尔伯特问题之一。


\subsection{数论}\label{Number_theory}

黎曼在现代\href{https://en.wikipedia.org/wiki/Analytic_number_theory}{解析数论}方面亦做出了若干闻名于世的贡献。他在数论方面发表的唯一一篇论文——\href{https://en.wikipedia.org/wiki/On_the_Number_of_Primes_Less_Than_a_Given_Magnitude}{一篇短文}中，研究了如今以他名字命名的\href{https://en.wikipedia.org/wiki/Riemann_zeta_function}{ζ 函数}，确立了它在理解\href{https://en.wikipedia.org/wiki/Prime_number}{素数}分布方面的重要性。\href{https://en.wikipedia.org/wiki/Riemann_hypothesis}{黎曼猜想}就是他对该函数性质所做的若干猜想之一。

黎曼的工作中还有许多有趣的进展。他证明了 ζ 函数的函数方程（其实\href{https://en.wikipedia.org/wiki/Leonhard_Euler}{莱昂哈德·欧拉}早已知晓），方程背后正是一个 theta 函数。通过把这个近似函数对实部为 1/2 的临界线上的非平凡零点求和，他给出了 {{\(\pi(x)\)}{}} 的精确"显式公式"。

黎曼当时已经知晓\href{https://en.wikipedia.org/wiki/Pafnuty_Chebyshev}{帕夫努蒂·切比雪夫（Chebyshev）}关于\href{https://en.wikipedia.org/wiki/Prime_number_theorem}{素数定理}的工作。切比雪夫曾于 1852 年访问过狄利克雷。


\subsection{著作}\label{Writings}

黎曼的主要著作包括：

\begin{itemize}
\tightlist
\item
  {1851 -- \emph{Grundlagen für eine allgemeine Theorie der Functionen einer veränderlichen complexen Grösse}（一个复变量函数一般理论的基础），博士论文，哥廷根，1851 年。}
\item
  {1857 -- \emph{\href{https://en.wikipedia.org/wiki/List_of_important_publications_in_mathematics\#Theorie_der_Abelschen_Functionen}{Theorie der Abelschen Functionen}}（阿贝尔函数论），\emph{Journal für die reine und angewandte Mathematik}（《纯粹与应用数学杂志》），第 54 卷，第 101--155 页。}
\item
  {1859 -- \emph{Über die Anzahl der Primzahlen unter einer gegebenen Größe}（论小于给定数值的素数个数），载\emph{《普鲁士科学院月刊》}（柏林，1859 年 11 月，第 671 页起），其中提出了黎曼猜想。\emph{Über die Anzahl der Primzahlen unter einer gegebenen Grösse.} (Wikisource)；\href{http://www.claymath.org/sites/default/files/riemann1859.pdf}{手稿影印本}（克莱数学研究所提供）。}
\item
  {1861 -- \emph{\href{https://www.maths.tcd.ie/pub/HistMath/People/Riemann/Paris/}{Commentatio mathematica, qua respondere tentatur quaestioni ab Illma Academia Parisiensi propositae}}（提交给巴黎科学院的有奖征文论文）。}
\item
  {1867 -- \emph{Über die Darstellbarkeit einer Function durch eine trigonometrische Reihe}（论函数可由三角级数表示），载哥廷根皇家科学学会论文集第 13 卷。}
\item
  {1868 -- \href{http://resolver.sub.uni-goettingen.de/purl?GDZPPN002019213}{\emph{Über die Hypothesen, welche der Geometrie zugrunde liegen}}（论作为几何学基础的假设），哥廷根皇家科学学会论文集，1868 年。\href{https://www.emis.de/classics/Riemann/WKCGeom.pdf}{英译本（EMIS, pdf）}\emph{On the hypotheses which lie at the foundation of geometry}，由\href{https://en.wikipedia.org/wiki/William_Kingdon_Clifford}{W.K. Clifford} 翻译，载 \emph{Nature} 第 8 卷（1873 年），第 183 页 -- 重印于克利福德的《数学论文集》（伦敦 MacMillan 1882；纽约 Chelsea 1968）。亦见 Ewald, William B. (主编), 1996, \emph{From Kant to Hilbert: A Source Book in the Foundations of Mathematics}（《从康德到希尔伯特：数学基础原始文献》），共 2 卷，牛津大学出版社：第 652--661 页。}
\item
  {1876 -- \emph{Bernhard Riemann\textquotesingle s Gesammelte Mathematische Werke und wissenschaftlicher Nachlass}（《黎曼数学全集与遗稿》），由海因里希·韦伯（Heinrich Weber）在理查德·戴德金协助下编辑，莱比锡，B. G. Teubner，1876 年；1892 年第二版；1953 年 Dover 重印（含马克斯·诺特与威廉·维尔廷格的贡献，Teubner 1902）。后续版本：\emph{The collected Works of Bernhard Riemann: The Complete German Texts}，由 Heinrich Weber、Richard Dedekind、M. Noether、Wilhelm Wirtinger、Hans Lewy 主编，纽约 Mineola：Dover Publications，1953、1981、2017 年。}
\item
  {1876 -- \emph{Schwere, Elektrizität und Magnetismus}（重力、电学与磁学），汉诺威：Karl Hattendorff。}
\item
  {1882 -- \emph{Vorlesungen über Partielle Differentialgleichungen}（偏微分方程讲义），第 3 版。布伦瑞克 1882 年。}
\item
  {1901 -- \emph{Die partiellen Differential-Gleichungen der mathematischen Physik nach Riemann\textquotesingle s Vorlesungen}（依黎曼讲座整理的数学物理偏微分方程）。维基共享资源上有 PDF。archive.org 上：Riemann, Bernhard (1901)；Weber, Heinrich Martin (主编)；\href{https://archive.org/details/diepartiellendi02webegoog}{"Die partiellen differential-gleichungen der mathematischen physik nach Riemann\textquotesingle s Vorlesungen"}，Friedrich Vieweg und Sohn 出版。}
\item
  {2004 -- Riemann, Bernhard (2004), \emph{Collected papers}（论文集），Kendrick Press, Heber City, UT，\href{https://en.wikipedia.org/wiki/ISBN_(identifier)}{ISBN}{~}\href{https://en.wikipedia.org/wiki/Special:BookSources/978-0-9740427-2-5}{978-0-9740427-2-5}，\href{https://en.wikipedia.org/wiki/MR_(identifier)}{MR}{~}\href{https://mathscinet.ams.org/mathscinet-getitem?mr=2121437}{2121437}{}}
\end{itemize}




延伸阅读：

\begin{itemize}
\tightlist
\item
  {\href{https://en.wikipedia.org/wiki/John_Derbyshire}{Derbyshire, John} (2003), \emph{\href{https://en.wikipedia.org/wiki/Prime_Obsession}{Prime Obsession: Bernhard Riemann and the Greatest Unsolved Problem in Mathematics}}（《素数之恋：黎曼与数学中最伟大的未解之谜》），华盛顿特区：John Henry Press，\href{https://en.wikipedia.org/wiki/ISBN_(identifier)}{ISBN}{~}\href{https://en.wikipedia.org/wiki/Special:BookSources/0-309-08549-7}{0-309-08549-7}.}
\item
  {Monastyrsky, Michael (1999), \emph{Riemann, Topology and Physics}（《黎曼、拓扑与物理》），波士顿：Birkhäuser，\href{https://en.wikipedia.org/wiki/ISBN_(identifier)}{ISBN}{~}\href{https://en.wikipedia.org/wiki/Special:BookSources/0-8176-3789-3}{0-8176-3789-3}.}
\item
  {Ji, Lizhen; Papadopoulos, Athanese; Yamada, Sumio, eds. (2017). \emph{From Riemann to Differential Geometry and Relativity}（《从黎曼到微分几何与相对论》），Springer。\href{https://en.wikipedia.org/wiki/ISBN_(identifier)}{ISBN}{~}\href{https://en.wikipedia.org/wiki/Special:BookSources/9783319600390}{9783319600390}.}
\end{itemize}



{{
\includegraphics[width=0.35417in,height=0.41667in]{40px-Wikiquote-logo.svg.png}}}

维基语录上有与\emph{\textbf{\href{https://en.wikiquote.org/wiki/Special:Search/Bernhard\%20Riemann}{波恩哈德·黎曼}}}相关的引言。


{\href{https://en.wikipedia.org/wiki/File:Commons-logo.svg}{
\includegraphics[width=0.3125in,height=0.41667in]{40px-Commons-logo.svg.png}}}

维基共享资源上有与\href{https://commons.wikimedia.org/wiki/Georg\%20Friedrich\%20Bernhard\%20Riemann}{波恩哈德·黎曼}相关的多媒体资源。


{\href{https://en.wikipedia.org/wiki/File:Wikisource-logo.svg}{
\includegraphics[width=0.39583in,height=0.41667in]{40px-Wikisource-logo.svg.png}}}

\href{https://en.wikipedia.org/wiki/Wikisource}{维基文库}收录有 1905 年版《\href{https://en.wikipedia.org/wiki/New_International_Encyclopedia}{新国际百科全书（New International Encyclopedia）}》中的"\textbf{Riemann, Georg Friedrich Bernhard}"词条。

\begin{itemize}
\tightlist
\item
  {\href{https://mathgenealogy.org/id.php?id=18232}{波恩哈德·黎曼} 在\href{https://en.wikipedia.org/wiki/Mathematics_Genealogy_Project}{数学家谱项目（Mathematics Genealogy Project）}的页面}
\item
  {\href{http://www.maths.tcd.ie/pub/HistMath/People/Riemann/Papers.html}{黎曼数学论文（The Mathematical Papers of Georg Friedrich Bernhard Riemann）}}
\item
  {\href{http://www.emis.de/classics/Riemann/}{emis.de 上黎曼的著作}}
\item
  {O\textquotesingle Connor, John J.; \href{https://en.wikipedia.org/wiki/Edmund_F._Robertson}{Robertson, Edmund F.}, \href{https://mathshistory.st-andrews.ac.uk/Biographies/Riemann.html}{"Bernhard Riemann"}, \emph{\href{https://en.wikipedia.org/wiki/MacTutor_History_of_Mathematics_Archive}{MacTutor 数学史档案库}}, \href{https://en.wikipedia.org/wiki/University_of_St_Andrews}{圣安德鲁斯大学}}
\item
  {\href{https://web.archive.org/web/20050903035028/http://www.fh-lueneburg.de/u1/gym03/englpage/chronik/riemann/riemann.htm}{波恩哈德·黎曼 -- 最重要的数学家之一}}
\item
  {\href{https://web.archive.org/web/20160318034045/http://www.maths.tcd.ie/pub/HistMath/People/Riemann/Geom/}{波恩哈德·黎曼的就职讲演}}
\item
  {\href{https://en.wikipedia.org/wiki/Eric_W._Weisstein}{Weisstein, Eric Wolfgang} (主编)，\href{http://scienceworld.wolfram.com/biography/Riemann.html}{"Riemann, Bernhard (1826--1866)"}，\emph{\href{https://en.wikipedia.org/wiki/ScienceWorld}{ScienceWorld}}.}
\item
  {\href{https://www.maths.tcd.ie/pub/HistMath/People/Riemann/Leben/Leben.pdf}{Richard Dedekind (1892)，由 D. R. Wilkins 转录的黎曼传记}}
\item
  {\href{https://en.wikipedia.org/wiki/David_E._Rowe}{David E. Rowe}: \href{https://www.ams.org/journals/notices/202602/noti3294/noti3294.html}{"Did Gauss Mentor Riemann?"}（"高斯指导过黎曼吗？"），\emph{Notices of AMS}，第 73 卷第 2 期（2026 年 2 月）。}
\end{itemize}




{}
